\chapter{Topological Endgame}

Finite extinction leaves a finite collection of standard components at the
last regular time.  This chapter reverses the finitely many preceding
surgeries, expresses the initial manifold as a connected sum of standard
factors, and then uses simple connectivity to eliminate every factor except
$S^3$.

\section{Smooth Connected Sums and Prime Factors}

\begin{definition}[Smooth connected sum]
\label{def:smooth-connected-sum}
\dcref{MT:ch15:15.3}
\source{morgan-tian}
\notready
Let $M$ and $N$ be connected smooth $3$-manifolds without boundary.  Given
smoothly embedded closed $3$-balls and a diffeomorphism of their boundary
spheres, their \emph{smooth connected sum} $M\mathbin{\#}N$ is obtained by
deleting the interiors of the balls and making that boundary identification.
In the oriented category the manifolds are oriented and the identification is
orientation reversing.  This orientation convention is part of the definition;
the isotopy-extension theorem for $3$-balls then identifies different choices
and makes finite connected sums associative up to oriented diffeomorphism.
\end{definition}

\begin{definition}[Spherical space form]
\label{def:spherical-space-form}
\dcref{MT:intro:0.1}
\source{morgan-tian}
\notready
A \emph{spherical space form} is a quotient $S^3/\Gamma$, where $S^3$ has
its round metric and $\Gamma$ is a finite group acting freely by isometries.
\end{definition}

\begin{definition}[The two sphere bundles over the circle]
\label{def:sphere-bundles-over-circle}
\dcref{MT:ch15:15.3}
\source{morgan-tian}
\notready
Write
\[
 E_+=S^2\times S^1
 \qquad\text{and}\qquad
 E_-=S^1\mathbin{\widetilde\times}S^2
\]
for, respectively, the orientable product bundle and the nonorientable
twisted smooth $S^2$-bundle over $S^1$.
\end{definition}

\begin{lemma}[The two smooth gluing classes of a two-sphere]
\label{lem:two-sphere-diffeomorphism-classes}
\dcref{MT:ch15:15.4}
\source{morgan-tian}
\uses{def:sphere-bundles-over-circle}
\notready
Every smooth diffeomorphism $\varphi\colon S^2\to S^2$ is smoothly isotopic
to the identity when it preserves orientation and to the antipodal map when it
reverses orientation.  Hence gluing the two sphere boundaries of
$S^2\times[0,1]$ produces exactly one of $E_+$ and $E_-$.  The resulting
closed three-manifold is orientable exactly in the $E_+$ case.
\end{lemma}

\begin{proof}
The mapping-class assertion and its smooth-isotopy regularity are the exact
two-sphere gluing interface used in the cited Morgan--Tian reconstruction.
Isotopic gluing maps give diffeomorphic mapping tori by extending the isotopy
over a collar.  The identity mapping torus is $E_+=S^2\times S^1$; the
orientation-reversing class gives $E_-$, and transport of a local
orientation once around the base gives the final orientability assertion.
\end{proof}

\begin{definition}[Simply connected space]
\label{def:simply-connected}
\dcref{MT:intro:0.1}
\source{morgan-tian}
\notready
A topological space is \emph{simply connected} when it is path-connected and
its fundamental group at one, equivalently every, basepoint is trivial.
\end{definition}

\begin{lemma}[Path-connected spaces are connected]
\label{lem:path-connected-implies-connected}
\dcref{MT:intro:0.1}
\source{morgan-tian}
\uses{def:simply-connected}
\notready
Every path-connected topological space is connected.  In particular, a space
which is simply connected in the sense of \ref{def:simply-connected} is
connected.
\end{lemma}

\begin{proof}
If $X=U\sqcup V$ were a separation and $x\in U$, $y\in V$, a path from $x$
to $y$ would have connected image contained in $U\sqcup V$, hence would lie
in one member of the separation, a contradiction.
\end{proof}

\begin{lemma}[A simply connected three-manifold is orientable]
\label{lem:simply-connected-three-manifold-orientable}
\dcref{MT:intro:0.1}
\source{morgan-tian}
\uses{def:simply-connected}
\notready
Every connected simply connected smooth $3$-manifold is orientable.
\end{lemma}

\begin{proof}
A connected simply connected manifold has no nontrivial connected double
cover.  Its orientation double cover is therefore trivial, so a consistent
choice of local orientation exists globally.
\end{proof}

\begin{theorem}[Hamilton roundness for positive Ricci curvature]
\label{thm:positive-ricci-three-manifold-roundness}
\dcref{MT:ch4:4.23}
\source{morgan-tian}
\uses{def:spherical-space-form}
\notready
Let $M$ be a compact connected smooth $3$-manifold without boundary and let
$g$ be a smooth Riemannian metric with $\operatorname{Ric}_g>0$ everywhere.
Then $M$ admits a metric of constant positive sectional curvature and is
diffeomorphic to a spherical space form.
\end{theorem}

\begin{proof}
Start the maximal smooth Ricci flow $(M,g(t))$, $0\le t<T$, at $g$.  The
positive-Ricci maximum principle in MT 4.23 preserves positivity, forces
$T<\infty$, and says that after the standard curvature rescaling the metrics
become round as $t\uparrow T$.  The underlying smooth manifold is fixed along
this flow, so the limiting round metric (transported through the flow
diffeomorphisms) is a constant-positive-curvature metric on $M$.  A compact
connected constant-positive-curvature $3$-manifold is $S^3/\Gamma$ for a finite
free isometric action, which is exactly
\ref{def:spherical-space-form}.  The convergence and space-form
classification are imported at the stated MT 4.23 interface.
\end{proof}

\section{Reversing One Surgery}

\begin{lemma}[Whole disappearing components]
\label{lem:whole-disappearing-component-classification}
\dcref{MT:ch4:4.23,MT:ch15:15.3,MT:app:A.25}
\source{morgan-tian}
\uses{def:controlled-ricci-flow-with-surgery,
  def:strong-canonical-neighborhood-control,
  def:spherical-space-form,
  def:sphere-bundles-over-circle,
  thm:positive-ricci-three-manifold-roundness,
  thm:canonical-neck-cap-topology}
\notready
Let $({\mathcal M},G)$ be a controlled Ricci flow with surgery satisfying
Assumptions (1)--(7) in the definition of controlled surgery.  Let $T$ be a
singular time and let $t^-<T$ be regular and sufficiently close to $T$.  Write
$C_T\subset M_T$ for the continuing region, let
$C_{t^-}\subset M_{t^-}$ be its backward time-flow image, and put
$D_{t^-}=M_{t^-}\setminus C_{t^-}$.  If a connected component of $D_{t^-}$ is
also a whole component of $M_{t^-}$, then it is diffeomorphic
to one of
\[
 S^3/\Gamma,\qquad \mathbb{RP}^3\mathbin{\#}\mathbb{RP}^3,
 \qquad E_+,
 \qquad E_-.
\]
\end{lemma}

\begin{proof}
By the controlled-surgery definition, every point of the component has a
strong canonical neighborhood immediately before it disappears.  If one of
these neighborhoods is a compact positively curved component, it is the
whole connected component and has positive Ricci curvature; it is a
spherical space form by
\ref{thm:positive-ricci-three-manifold-roundness}.  The epsilon-round
component alternative is $C^2$-close, after rescaling, to a round metric.
Since epsilon is chosen sufficiently small, its Ricci curvature is positive,
so the same theorem applies.

In the remaining cases every point is a neck center or belongs to a cap core.
Compactness gives a finite neck-and-cap cover.  The compact closed
alternatives in \ref{thm:canonical-neck-cap-topology} are $S^3$,
$\mathbb{RP}^3$, $\mathbb{RP}^3\#\mathbb{RP}^3$, or an $S^2$-bundle over
$S^1$.  The first two are spherical space forms, and the last alternative is
$E_+$ or $E_-$ according to its orientability.
\end{proof}

\begin{definition}[Proper disappearing component]
\label{def:proper-disappearing-component}
\dcref{MT:ch15:15.3}
\source{morgan-tian}
\uses{def:controlled-ricci-flow-with-surgery}
\notready
Let $({\mathcal M},G)$ satisfy Assumptions (1)--(7), let $T$ be a singular
time, and let $t^-<T$ be regular and sufficiently close to $T$.  Write
$C_{t^-}$ for the backward image of the continuing region and
$D_{t^-}=M_{t^-}\setminus C_{t^-}$.  A connected component $D$ of
$D_{t^-}$ is \emph{proper} when it is not a connected component of the ambient
slice $M_{t^-}$.  Its frontier always means the ambient frontier
$\operatorname{Fr}_{M_{t^-}}D$; no spherical-frontier, compactness, or
canonical-neighborhood conclusion is included in this definition.
\end{definition}

\begin{lemma}[Actual disappearing components have saturated neck frontier]
\label{lem:actual-disappearing-component-frontier-saturation}
\dcref{MT:ch15:15.2,MT:ch15:15.3}
\source{morgan-tian}
\uses{def:controlled-ricci-flow-with-surgery,
  def:proper-disappearing-component,def:embedded-horn-frontier,
  def:canonical-neck-cap-regions}
\notready
Let $({\mathcal M},G)$ satisfy Assumptions (1)--(7), let $t$ be a singular
time, and let $t^-<t$ be regular and sufficiently close to $t$.  If a connected
component $D$ of the actual disappearing region
$D_{t^-}=M_{t^-}\setminus C_{t^-}$ is not a whole connected component of
$M_{t^-}$ (where the presurgery slice is compact, smooth, and boundaryless),
then $D$ has nonempty frontier and $\overline D$ is a
compact codimension-zero smooth submanifold,
\[
 \operatorname{Fr}_{M_{t^-}}D=S_1\sqcup\cdots\sqcup S_m
\]
is a finite union of pairwise disjoint embedded $2$-spheres, and every
$S_j$ is the backward image of a component of $\Sigma_t$ and hence the central
sphere of the corresponding strong neck in the maximal presurgery extension.
Every point of $D$ is either an $\epsilon$-neck center or belongs to the core
of a $(C,\epsilon)$-cap.
\end{lemma}

\begin{proof}
The definition of the continuing region, together with Assumption (4),
writes $C_t$ as a compact codimension-zero submanifold with finite spherical
boundary and identifies it by backward flow with the compact
codimension-zero submanifold $C_{t^-}\subset M_{t^-}$.  Hence
$D_{t^-}=M_{t^-}\setminus C_{t^-}$ is open, each of its connected components
is maximal among connected subsets of $D_{t^-}$, and the frontier of any such
component is a subcollection of the finite smooth boundary
$\partial C_{t^-}$.  A component has empty frontier exactly when it is open
and closed in its ambient connected component; the hypothesis that $D$ is not
a whole component therefore makes this subcollection nonempty.  Its closure
is the union of $D$ with these boundary components and is consequently a
compact codimension-zero smooth submanifold.

Assumption (5) identifies every component of $\partial C_{t^-}$, by backward
flow, with the central sphere of a strong neck in the maximal presurgery
extension.  Assumption (6) gives a strong canonical neighborhood at every
point of $D$.  The compact positively curved and epsilon-round alternatives
are entire connected components of the slice, so neither can occur in the
present frontier-bearing component.  Thus every point of $D$ is a neck center
or lies in a canonical-cap core.  These are exactly the canonical-cover
clauses asserted above.
\end{proof}

\begin{lemma}[Fiber saturation from a spherical frontier]
\label{lem:fiber-saturation-from-spherical-frontier}
\dcref{MT:app:A.20,MT:app:A.21}
\source{morgan-tian}
\lean{PoincareConjecture.Topology.FiberSaturation.fiber_saturation_from_spherical_frontier}
\leanok
Let $J\subset\mathbb R$ be a nonempty open interval and let
$U\subset S^2\times J$ be connected and open.  Suppose that
$\operatorname{Fr}_{S^2\times J}(U)$ is a union of complete fibers
$S^2\times\{s\}$.  Then there is an open interval $A\subset J$ such that
$U=S^2\times A$.  If, in addition, $\overline U$ is compact in
$S^2\times J$ and the frontier is nonempty, then there are $a<b$ in the
interior of $J$ such that
\[
 U=S^2\times(a,b),\qquad
 \operatorname{Fr}_{S^2\times J}(U)
 =S^2\times\{a\}\,\sqcup\,S^2\times\{b\}.
\]
\end{lemma}

\begin{proof}
For a parameter $s$ whose fiber is not in the frontier, the sets
$U\cap(S^2\times\{s\})$ and
$(S^2\times J\setminus\overline U)\cap(S^2\times\{s\})$ are disjoint open
subsets of the connected sphere $S^2$ and cover it.  Thus that fiber is
either wholly contained in $U$ or wholly outside $\overline U$; frontier
fibers are wholly in the frontier by hypothesis.  Hence
$U=S^2\times A$ for an open subset $A\subset J$.  Connectedness of $U$
makes $A$ an interval.  Since $J$ is open, compactness of $\overline U$ and
nonempty frontier make its endpoints finite and interior to $J$, and the
frontier is exactly the two endpoint fibers.  (Without the openness of $J$, a
relative-open set could end at an endpoint of $J$ and have only one relative
frontier fiber.)  This elementary saturation argument is the
relative step used in the neck-chain classification of Morgan--Tian
Appendices~A.20--A.21.
\end{proof}

\begin{lemma}[Saturation on a closed product collar]
\label{lem:closed-collar-fiber-saturation}
\dcref{MT:app:A.20,MT:app:A.21}
\source{morgan-tian}
\uses{lem:fiber-saturation-from-spherical-frontier}
\notready
Let $E=S^2\times[a,b]$ be a closed product collar in a smooth
three-manifold $Y$, with endpoint fibers $F_a$ and $F_b$.  Let $D\subset Y$ be
connected and open, with compact smooth codimension-zero closure and
\(\partial\overline D=\operatorname{Fr}_Y(D)\) a finite union of displayed
cross-sectional spheres.  For every connected component $V$ of
$D\cap(S^2\times(a,b))$, there is an open interval $A\subset(a,b)$ such that
\[
 V=S^2\times A.
\]
Moreover, if $D$ meets an endpoint fiber, it contains that whole fiber; if an
endpoint fiber meets $\operatorname{Fr}_Y(D)$, that whole fiber is a frontier
component.  Thus $D$ crosses each product collar by whole fibers, and can stop
there only at a displayed cross-section.
\end{lemma}

\begin{proof}
For a component $V$ as in the statement, any point of its frontier in the
open collar is a frontier point of $D$ in $Y$: two distinct components of an
open set cannot meet without the intervening complement, and that complement
is part of the ambient frontier.  The displayed-frontier hypothesis therefore
makes $\operatorname{Fr}_{S^2\times(a,b)}(V)$ a union of complete fibers.
Apply the first clause of
\ref{lem:fiber-saturation-from-spherical-frontier} with
$J=(a,b)$ to obtain $V=S^2\times A$.  At an endpoint, a partial intersection
with $F_a$ or $F_b$ would have a boundary inside that sphere; a small product
collar would then produce a frontier point which is not a complete displayed
fiber.  Hence an endpoint fiber is either wholly in $D$, wholly outside
$\overline D$, or wholly in the frontier.  This proves all assertions and is
the endpoint-safe collar form of the relative saturation step in
Appendices~A.20--A.21.
\end{proof}

\begin{lemma}[Boundary pairing in a fibered neck region]
\label{lem:relative-fiber-boundary-pairing}
\dcref{MT:app:A.20,MT:app:A.21,MT:app:A.24}
\source{morgan-tian}
\uses{lem:fiber-saturation-from-spherical-frontier}
\notready
Let $E$ be either an open product $S^2\times J$ with $J$ an open interval,
or a closed $S^2$-bundle over $S^1$ presented as

\[
 q:S^2\times\mathbb R\longrightarrow E,
 \qquad q(y,s+L)=q(\phi(y),s)
\]

for some $L>0$ and diffeomorphism $\phi:S^2\to S^2$.  Let $D\subset E$ be
connected and open, with compact smooth codimension-zero closure satisfying
\(\partial\overline D=\operatorname{Fr}_E(D)\ne\varnothing\), and suppose
that this frontier is a union of complete fibers.  Then

\[
 \overline D\cong S^2\times[0,1],
 \qquad \operatorname{Fr}_E(D)=S_0\sqcup S_1.
\]
\end{lemma}

\begin{proof}
In the product case apply
\ref{lem:fiber-saturation-from-spherical-frontier} directly.  In the bundle
case, every component of $q^{-1}(D)$ has frontier a union of complete lifted
fibers, so the first (noncompact) part of the saturation argument gives one
component $S^2\times(a,b)$.  Distinct translates of this interval are
disjoint, since they are distinct components of the preimage.  Hence
$b-a\le L$; if $b-a>L$ they overlap, while if $b-a=L$ the quotient closure is
locally the whole mapping-torus neighborhood at the identified endpoint and
has no boundary there, contradicting
\(\partial\overline D=\operatorname{Fr}_E(D)\ne\varnothing\).  Thus
$b-a<L$, the quotient is injective on $S^2\times[a,b]$, and its two endpoint
fibers are distinct boundary components.  The product collar gives the stated
diffeomorphism.  This is the endpoint-pairing argument implicit in the
relative use of Appendices~A.20--A.24.
\end{proof}

\begin{lemma}[Relative classification of a proper all-neck region]
\label{lem:proper-all-neck-relative-fibration}
\dcref{MT:app:A.20,MT:app:A.21,MT:app:A.24}
\source{morgan-tian}
\uses{def:proper-disappearing-component,
  lem:actual-disappearing-component-frontier-saturation,
  prop:overlapping-epsilon-necks,def:balanced-neck-chain,
  lem:neck-chain-product,lem:infinite-neck-chain-frontier,
  lem:all-neck-component-tube-or-bundle,
  lem:fiber-saturation-from-spherical-frontier,
  lem:relative-fiber-boundary-pairing,
  lem:neck-separation-type-locally-constant,
  thm:separating-neck-centers-form-tube,
  lem:nonseparating-neck-chain-or-bundle}
\notready
Let $D$ be a proper connected component of a disappearing region in a compact
boundaryless presurgery $3$-manifold.  Assume every point of $D$ is an
$\epsilon$-neck center, every point of its nonempty frontier is a central
neck sphere, and $\overline D$ is a compact smooth codimension-zero
submanifold with
\[
 \partial\overline D=\operatorname{Fr}D\ne\varnothing
\]
a finite union of those spheres.  Then
\[
 \overline D\cong S^2\times[0,1],\qquad
 \operatorname{Fr}D=S_0\sqcup S_1.
\]
In particular, the closed $S^2$-bundle alternative in Appendix~A.20 cannot
be used to replace this relative conclusion: Appendix~A.21 only gives
containment in an ambient neck-covered component.
\end{lemma}

\begin{proof}
The separation type of the neck centers is locally constant by
\ref{lem:neck-separation-type-locally-constant}.  In the separating case,
\ref{thm:separating-neck-centers-form-tube} supplies a product chart for a
balanced chain containing $D$.  The overlap/product construction straightens
the finitely many frontier central spheres simultaneously to fibers of that
chart.  The boundary-pairing lemma in the open-product case gives the claimed
two fibers and product closure.

In the nonseparating case, use
\ref{lem:nonseparating-neck-chain-or-bundle}.  Its tube alternative is the
same open-product case.  In its fibration alternative,
\ref{lem:all-neck-component-tube-or-bundle} identifies the ambient component
with a mapping torus of the fiber-preserving neck cover.  After straightening
the lifted frontier spheres to fibers, the boundary-pairing lemma in the
mapping-torus case applies.  Its endpoint argument uses the actual-frontier
identity
\(\partial\overline D=\operatorname{Fr}D\ne\varnothing\) from
\ref{lem:actual-disappearing-component-frontier-saturation}, so an identified
endpoint (which would make the quotient closure locally boundaryless) is
excluded.  This is the relative saturation step implicit in the source proof
of Proposition~15.3; A.24 is used only for cap gluing.
\end{proof}

\begin{lemma}[Exposed-end successors share a middle overlap]
\label{lem:exposed-end-successor-overlap}
\dcref{MT:app:A.11,MT:app:A.17,MT:app:A.21}
\source{morgan-tian}
\uses{def:balanced-neck-chain,prop:overlapping-epsilon-necks,
  lem:epsilon-neck-metric-control,lem:finite-neck-chain-frontier-dichotomy,
  lem:capped-chain-exposed-frontier}
\notready
For sufficiently small $\epsilon$, let $N_a,\ldots,N_b$ be a balanced chain
with only the positive end exposed.  Every neck centered at a frontier point
which satisfies the extension hypotheses of the finite-chain frontier
dichotomy (or the capped-chain exposed-frontier lemma) meets the middle
$0.9\epsilon^{-1}$ of $N_b$.  Any two such successor necks have a common
cross-sectional sphere in the middle overlap, isotopic in both necks to their
central spheres.  Consequently they determine one maximal product collar and
one successor incidence, rather than two branches.
\end{lemma}

\begin{proof}
The exposed-end hypotheses place each successor center in the positive end
overlap of $N_b$.  Neck metric control bounds the intrinsic displacement from
that center to the $0.9\epsilon^{-1}$ middle of $N_b$ by the neck scale, while
the overlap proposition compares the two scales by $1\pm\alpha$ and aligns
the axial directions.  For small enough $\epsilon$ (with the fixed
$\alpha\le10^{-2}$), the two axial intervals therefore overlap through a
common middle slice.  The middle-overlap clause of
\ref{prop:overlapping-epsilon-necks} supplies a cross-sectional sphere there
and its isotopies in both charts.  The product collar between that sphere and
the exposed end is independent of the chosen center, so all successor necks
represent the same maximal collar.  This is the quantitative overlap step in
Appendix~A.11 used by the chain extensions in A.17 and A.21.
\end{proof}

\begin{lemma}[Maximal-overlap incidence has no branching]
\label{lem:maximal-overlap-incidence-no-branching}
\dcref{MT:app:A.17,MT:app:A.21,MT:app:A.23,MT:app:A.25}
\source{morgan-tian}
\uses{def:canonical-neck-cap-regions,def:balanced-neck-chain,
  prop:overlapping-epsilon-necks,lem:neck-chain-product,
  lem:balanced-neck-chain-extension,lem:finite-neck-chain-frontier-dichotomy,
  lem:capped-chain-exposed-frontier,lem:finite-cap-chain-frontier-dichotomy,
  lem:infinite-neck-chain-frontier,lem:adjacent-canonical-caps-form-component,
  lem:exposed-end-successor-overlap}
\notready
Let $Y$ be a connected union of sufficiently small canonical necks and cap
cores.  Form its incidence graph from \emph{maximal} connected product collars:
two overlapping neck pieces are one edge whenever their middle overlap contains
a common cross-sectional sphere isotopic in both necks, and cap cores are
vertices.  Assume the finite-chain frontier extensions and the cap frontier
localization clauses of Appendices~A.17 and A.21 hold, and that an infinite
chain has no frontier.  Then every neck vertex has degree at most two, every
cap vertex is terminal, and the connected incidence graph is a (possibly
degenerate) finite interval, a half-line, a line, or a circle.
\end{lemma}

\begin{proof}
Each neck has exactly two end fibers, so a maximal product collar can have at
most one predecessor and one successor.  Suppose two distinct successor edges
leave the same exposed end of a balanced chain.  By the finite-neck frontier
  dichotomy and the capped-chain localization, both successor centers lie in the
  same connected exposed-end region.  Apply
  \ref{lem:exposed-end-successor-overlap}: the successor necks meet a common
  middle slice and their cross-sectional spheres are isotopic in both charts.
  Projection along the axial lines identifies both collars with the same
  $S^2\times I$ product.  Maximality
of the product-collar equivalence relation makes the two supposed successors
one edge, a contradiction.  The same argument at the negative end handles
predecessors.  A second cap at an exposed end either is nested and would have
been selected by the finite cap dichotomy, or satisfies the
$C_0\nsubseteq C'$ alternative of the adjacent-cap lemma and closes the two
ends into one component; it cannot create a third incident edge.

The overlap construction makes the graph connected because $Y$ is connected.
Thus it is a connected graph with all degrees at most two.  Such a graph is a
finite interval, a ray, a bi-infinite line, or a circle (with the one-vertex
case included as a degenerate interval).  The no-frontier lemma rules out an
additional vertex at an infinite end, and the cap cores are terminal vertices
by definition.  This is the explicit no-branching reduction underlying the
six containment alternatives in Appendix~A.21 and the global alternatives of
Appendix~A.25.
\end{proof}

\begin{lemma}[Ordered incidence supplied by the cap/neck classification]
\label{lem:canonical-cover-incidence-structure}
\dcref{MT:app:A.21,MT:app:A.25}
\source{morgan-tian}
\uses{thm:cap-and-neck-cover-classification,
  lem:maximal-overlap-incidence-no-branching,
  def:canonical-neck-cap-regions,lem:neck-chain-product,
  lem:infinite-neck-chain-frontier}
\notready
Every containing component furnished by
\ref{thm:cap-and-neck-cover-classification} has the following additional
source structure.  Its cap/neck incidence graph is a finite interval, a
half-line, a line, or a circle; each edge is a product collar
$S^2\times I$, each cap is a terminal vertex, and all central neck spheres
are cross-sectional fibers in the corresponding edge.  The incidence graph
has no branching.
\end{lemma}

\begin{proof}
The six alternatives in
\ref{thm:cap-and-neck-cover-classification} are respectively two terminal
caps, a doubly capped tube, one cap, a singly capped tube, an uncapped tube,
or a closed fibration.  Apply
\ref{lem:maximal-overlap-incidence-no-branching} to the maximal product-collar
decomposition.  Its degree bound gives a single axial order through every
successive neck overlap; its cap clause makes the cap cores terminal, and its
infinite-end clause removes a hidden frontier at an infinite edge.  The return
case closes the order into the circle of the fibration alternative.  Hence the
incidence graph has exactly the listed one-dimensional forms, with no branch
vertex.
\end{proof}

\begin{lemma}[Ordered incidence of a relative capped cover]
\label{lem:ordered-relative-capped-cover}
\dcref{MT:app:A.21,MT:app:A.23,MT:app:A.24,MT:app:A.25}
\source{morgan-tian}
\uses{lem:canonical-cover-incidence-structure,
  def:canonical-neck-cap-regions,prop:overlapping-epsilon-necks,
  lem:neck-chain-product,lem:infinite-neck-chain-frontier,
  lem:fiber-saturation-from-spherical-frontier,
  lem:closed-collar-fiber-saturation}
\notready
Let $Y$ be a connected canonical cap/neck component with the ordered
incidence structure supplied by Appendix~A.21: after cutting along its central
neck spheres, the incidence graph is a finite interval, a half-line, a line,
or a circle; every edge is an $S^2\times I$ product collar and every cap is a
terminal vertex with core diffeomorphic to a $3$-ball or to
$\mathbb{RP}^3\setminus\operatorname{int}(B^3)$.  Let $D\subset Y$ be a
proper connected open set with compact smooth codimension-zero closure, whose
nonempty frontier is a finite union of the displayed cross-sectional spheres
and satisfies
\[
 \partial\overline D=\operatorname{Fr}_Y(D)\ne\varnothing.
\]
If $D$ meets a cap core, then $\operatorname{Fr}_Y(D)$ has exactly one
component and

\[
 \overline D\cong B^3
 \quad\text{or}\quad
 \overline D\cong\mathbb{RP}^3\setminus\operatorname{int}(B^3).
\]
\end{lemma}

\begin{proof}
Cut the ordered cover along all its displayed central spheres.  The maximal
overlap lemma and the neck-chain product lemma identify the result with a graph
of product collars whose underlying graph is the stated interval, ray, line,
or circle; the cap cores are precisely the terminal vertices.  In an edge
collar $E=S^2\times[a,b]$, apply
\ref{lem:closed-collar-fiber-saturation} to each component of
$D\cap(S^2\times(a,b))$.  Thus the intersection is a union of whole fibers,
and it can stop at an endpoint only along the complete endpoint cross-section.
Because the frontier of $D$ contains no cap-interior points, meeting a cap core
forces $D$ to contain that entire connected cap core; otherwise the boundary of
the partial cap intersection would be a non-displayed surface.  Hence $D$
determines a connected subgraph of the incidence graph, and its boundary is
exactly the set of edge cross-sections at which this saturated subgraph stops.

A cap is a terminal vertex, so a proper connected subgraph containing a cap
has one outgoing edge: in a finite interval, containing both terminal caps
would force the whole interval; in a ray or line the cap end is the unique
terminal direction; and a circle has no cap vertex.  The whole-interval case
is excluded by properness, while an infinite ray/line end contributes no
additional frontier by \ref{lem:infinite-neck-chain-frontier}.  Therefore the
relative subgraph has exactly one stopping edge, and its cross-sectional
sphere is the sole frontier component.  Gluing its product collar to the
terminal cap core preserves the two possible core topologies, proving the
displayed alternatives.  This is the relative, side-by-side form of the
three cap-incidence subcases in Appendices~A.21--A.24.
\end{proof}

\begin{lemma}[Capped proper regions have one spherical frontier]
\label{lem:proper-capped-region-classification}
\dcref{MT:ch15:15.3,MT:app:A.21,MT:app:A.22,MT:app:A.23,MT:app:A.24}
\source{morgan-tian}
\uses{def:proper-disappearing-component,
  lem:actual-disappearing-component-frontier-saturation,
  thm:cap-and-neck-cover-classification,
  prop:overlapping-epsilon-necks,lem:neck-chain-product,
  lem:ordered-relative-capped-cover,
  def:canonical-neck-cap-regions}
\notready
Let $D$ be a proper connected component of a disappearing region in a compact
boundaryless $3$-manifold.  Assume that every point of $D$ is an
$\epsilon$-neck center or lies in a $(C,\epsilon)$-cap core, that
$\overline D$ is a compact smooth codimension-zero submanifold with nonempty
finite spherical frontier, and every frontier component is a central neck
sphere.  If $D$ meets a cap core, then
\[
 \overline D\cong B^3
 \quad\text{or}\quad
 \overline D\cong\mathbb{RP}^3\setminus\operatorname{int}(B^3),
 \qquad
 \operatorname{Fr}D\cong S^2.
\]
\end{lemma}

\begin{proof}
The actual-frontier lemma supplies the compact smooth closure, the equality
\(\partial\overline D=\operatorname{Fr}D\), and the finite central-sphere
frontier.  The overlap/product isotopies straighten those finitely many
central spheres to the displayed cross-sectional fibers of the containing
ordered cover.  Apply the cap/neck cover theorem
\ref{thm:cap-and-neck-cover-classification} to the connected set $D$.
It places $D$ in a canonical component with the ordered incidence structure
of Appendix~A.21.  The relative ordered-cover lemma
\ref{lem:ordered-relative-capped-cover} applies to that component: its proof
keeps the side of each central sphere lying in $D$, uses the A.23 overlap
subcases to rule out branching, and uses the A.22/A.24 no-infinite-end and
cap-gluing clauses.  It gives one frontier sphere and attaches one product
collar to the cap core.  The two possible cap-core topologies are exactly
$B^3$ and $\mathbb{RP}^3\setminus\operatorname{int}(B^3)$, proving the
displayed alternatives.  This is the relative form of the capped case in
Morgan--Tian Proposition~15.3, with no containment-to-equality inference.
\end{proof}

\begin{lemma}[Relative-boundary classification for a proper disappearing region]
\label{lem:frontier-disappearing-component-classification}
\dcref{MT:ch15:15.3,MT:app:A.20,MT:app:A.21,MT:app:A.24}
\source{morgan-tian}
\uses{def:proper-disappearing-component,
  lem:actual-disappearing-component-frontier-saturation,
  def:canonical-neck-cap-regions,
  def:embedded-horn-frontier,
  prop:overlapping-epsilon-necks,
  lem:neck-chain-product,
  lem:balanced-neck-chain-extension,
  lem:neck-separation-type-locally-constant,
  thm:separating-neck-centers-form-tube,
  lem:nonseparating-neck-chain-or-bundle,
  lem:all-neck-component-tube-or-bundle,
  lem:infinite-neck-chain-frontier,
  lem:fiber-saturation-from-spherical-frontier,
  lem:no-infinite-nested-canonical-caps,
  lem:adjacent-canonical-caps-form-component,
  thm:cap-and-neck-cover-classification,
  lem:isotopic-cap-gluing-sphere,
  lem:whole-disappearing-component-classification,
  lem:proper-all-neck-relative-fibration,
  lem:proper-capped-region-classification}
\notready
Let $({\mathcal M},G)$ satisfy Assumptions (1)--(7), let $T$ be a singular
time, and let $t^-<T$ be regular and sufficiently close to $T$.  If $D$ is a
proper connected component of $D_{t^-}$ in the sense of
\ref{def:proper-disappearing-component}, then the closure
$\overline D$ in $M_{t^-}$ is a compact codimension-zero smooth submanifold,
and there are $m\in\{1,2\}$ pairwise disjoint embedded spheres
$S_1,\ldots,S_m$ such that
\[
 \partial_{M_{t^-}}\overline D
 =\operatorname{Fr}_{M_{t^-}}D=S_1\sqcup\cdots\sqcup S_m,
\]
with each $S_j$ the backward image of a surgery central sphere.  The
source-exact neck/cap case analysis gives precisely one of
\[
 \overline D\cong S^2\times[0,1],\qquad
 \overline D\cong B^3,\qquad
 \overline D\cong\mathbb{RP}^3\setminus\operatorname{int}(B^3).
\]
The first alternative has two frontier spheres and each of the other two has
one.  No assertion is made here for a disappearing component with empty
frontier; that case is covered by
\ref{lem:whole-disappearing-component-classification}.
\end{lemma}

\begin{proof}
\uses{def:controlled-ricci-flow-with-surgery,
  def:proper-disappearing-component,
  lem:actual-disappearing-component-frontier-saturation,
  thm:cap-and-neck-cover-classification,
  thm:canonical-neck-cap-topology}
The preceding source-backed saturation lemma supplies compact closure, the
nonempty finite spherical frontier, and the cap/neck cover.  In particular, the
Appendix containment statement is not being used as a substitute for those
frontier facts.  The proof of Morgan--Tian Proposition~15.3 in
\texttt{surgery.tex}, lines 1544--1574, now applies directly to this proper
component.

If the component contains a cap core, apply
\ref{lem:proper-capped-region-classification}.  Its path-extension proof
excludes a second cap by the adjacent-cap component lemma and leaves one
exposed central sphere; the cap core and its product collar give exactly
$B^3$ or $\mathbb{RP}^3\setminus\operatorname{int}(B^3)$.

  If every point is a neck center, apply
  \ref{lem:proper-all-neck-relative-fibration}.  Its relative saturation
  argument is what excludes the invalid inference from containment in an
  $\epsilon$-fibration to equality with the closed bundle.  It gives exactly
  two frontier spheres and the $S^2\times[0,1]$ closure.

These source cases give precisely the three alternatives in the statement,
with the indicated one-versus-two frontier count.  The empty-frontier case is
the separate whole-component contract
\ref{lem:whole-disappearing-component-classification}; it is not folded into
this relative statement.
\end{proof}

\begin{lemma}[Connected-sum effect of capping frontier spheres]
\label{lem:local-capping-connected-sum-identities}
\dcref{MT:ch15:15.3}
\source{morgan-tian}
\uses{def:smooth-connected-sum,def:sphere-bundles-over-circle,
  lem:two-sphere-diffeomorphism-classes,
  lem:frontier-disappearing-component-classification,
  def:embedded-horn-frontier}
\notready
Let $X$ be a connected closed smooth $3$-manifold and let
$C\subset X$ be a compact codimension-zero smooth submanifold whose relative
boundary is the frontier-sphere union in
\ref{lem:frontier-disappearing-component-classification}.  Use collar
coordinates to identify each boundary sphere with the corresponding central
neck sphere, and let $Y$ be obtained by deleting $\operatorname{int}C$ and
capping every resulting boundary sphere by a smooth $3$-ball.
\begin{enumerate}
\item If $C\cong B^3$, then $X\cong Y$.
\item If $C\cong\mathbb{RP}^3\setminus\operatorname{int}(B^3)$, then
  $X\cong Y\#\mathbb{RP}^3$.
\item If $C\cong S^2\times[0,1]$ and its middle sphere separates $X$, then
  $Y=Y_1\sqcup Y_2$ and $X\cong Y_1\#Y_2$.
\item If $C\cong S^2\times[0,1]$ and its middle sphere does not separate $X$,
  then $Y$ is connected and $X\cong Y\#E_\varepsilon$ for some
  $\varepsilon\in\{+,-\}$.  If $X$ is orientable, then $\varepsilon=+$.
\end{enumerate}
\end{lemma}

\begin{proof}
Replacing a ball by a ball does not change the diffeomorphism type because
every diffeomorphism of $S^2$ extends across $B^3$.  Replacing a ball by
$\mathbb{RP}^3\setminus\operatorname{int}(B^3)$ is exactly the connected-sum
construction with $\mathbb{RP}^3$.

For the product region, cut $X$ along its middle sphere and cap the two copies
of that sphere.  If the cut separates, the two capped pieces are $Y_1$ and
$Y_2$, so reversing the cut is the definition of $Y_1\#Y_2$.  If the cut is
nonseparating, the capped complement is $Y$.  Reversing the operation removes
two balls from $Y$ and identifies their boundary spheres.  A regular
neighborhood of an arc joining the two balls isolates the mapping torus of
that boundary identification, so the result is $Y\#E_\varepsilon$.  The two
possibilities and their orientability are the nonirreducible prime alternatives
in \ref{lem:two-sphere-diffeomorphism-classes}; hence an orientable $X$ forces
$E_\varepsilon=E_+$.
\end{proof}

\begin{theorem}[Morgan--Tian one-surgery topological type]
\label{thm:morgan-tian-proposition-15-3}
\dcref{MT:ch15:15.3}
\source{morgan-tian}
\uses{def:controlled-ricci-flow-with-surgery,
  def:ricci-flow-with-cutoff,
  lem:actual-disappearing-component-frontier-saturation,
  lem:whole-disappearing-component-classification,
  lem:frontier-disappearing-component-classification,
  lem:local-capping-connected-sum-identities}
\notready
Let $({\mathcal M},G)$ be a controlled generalized Ricci flow satisfying
Assumptions (1)--(7) of Morgan--Tian and carrying the Kleiner--Lott
$(r,\delta)$-cutoff of \ref{def:ricci-flow-with-cutoff} on the relevant
finite interval.  If $T$ is a singular time and $t^-<T$ is regular and
sufficiently close to $T$, then, for each connected component $X$ of
$M_{t^-}$, there is a finite connected-sum expression whose summands are
connected components of $M_T$, smooth $S^2$-bundles over $S^1$, and closed
constant-positive-curvature three-manifolds.  Taking the disjoint union over
the components $X$ gives the corresponding global statement.  The number of
summands and operations is finite because the exposed region and its frontier
are finite.
\end{theorem}

\begin{proof}
Assumption (4) writes $M_T=C_T\cup_{\Sigma_T}E_T$ with $E_T$ a finite
disjoint union of three-balls and $\Sigma_T$ a finite union of two-spheres.
Backward flow identifies $C_T$ with $C_{t^-}\subset M_{t^-}$.  Consequently
$M_T$ is obtained from $M_{t^-}$ by deleting the interior of
$D_{t^-}=M_{t^-}\setminus C_{t^-}$ and capping every component of
$\partial C_{t^-}$ by a three-ball.

  By Assumption (6), every point of $D_{t^-}$ has the stated strong
  canonical neighborhood.  A component of $D_{t^-}$ which is a whole
  component of $M_{t^-}$ is therefore one of the alternatives in
  \ref{lem:whole-disappearing-component-classification}.  A component with
  nonempty relative frontier satisfies the maximality and central-neck
  hypotheses by
  \ref{lem:actual-disappearing-component-frontier-saturation}, so
  \ref{lem:frontier-disappearing-component-classification} applies;
it is a ball, a punctured $\mathbb{RP}^3$, or a product
$S^2\times[0,1]$ with precisely the indicated frontier spheres.

Apply \ref{lem:local-capping-connected-sum-identities} to each
frontier-bearing component.  A ball contributes nothing; a punctured
$\mathbb{RP}^3$ contributes the positive-curvature factor
$\mathbb{RP}^3$; and a product either joins two capped components by a
connected sum or contributes an $S^2$-bundle over $S^1$ in the
nonseparating case.  Whole disappearing components contribute the same
standard factors directly, with
$\mathbb{RP}^3\#\mathbb{RP}^3$ counted as two positive-curvature factors.
There are finitely many cases because the frontier and the time slices are
compact.  Combining them gives the asserted componentwise descriptions, and
hence the global disjoint-union and connected-sum description.
\end{proof}

\begin{lemma}[Reconstruction across one surgery time]
\label{lem:single-surgery-connected-sum-reconstruction}
\dcref{MT:ch15:15.3}
\source{morgan-tian}
\uses{def:controlled-ricci-flow-with-surgery,
  def:ricci-flow-with-cutoff,
  def:smooth-connected-sum,
  def:spherical-space-form,
  def:sphere-bundles-over-circle,
  lem:actual-disappearing-component-frontier-saturation,
  lem:whole-disappearing-component-classification,
  lem:frontier-disappearing-component-classification,
  lem:local-capping-connected-sum-identities}
\notready
Let $({\mathcal M},G)$ be a controlled generalized Ricci flow satisfying
Assumptions (1)--(7) and carrying the Kleiner--Lott cutoff of
\ref{def:ricci-flow-with-cutoff} on the relevant finite interval.  Let $T$ be a
singular time, and let $t^-<T$ be regular and sufficiently close to $T$.
Every component of $M_{t^-}$ is a finite connected sum of components of $M_T$
and factors
belonging to
\[
 \{S^3/\Gamma,\mathbb{RP}^3,E_+,E_-\}.
\]
Here a disappearing $\mathbb{RP}^3\#\mathbb{RP}^3$ component contributes two
$\mathbb{RP}^3$ factors.
\end{lemma}

\begin{proof}
  The frontier spheres are the finite disjoint union supplied by Assumption (4),
and the backward image $C_{t^-}$ is identified with the continuing region
$C_T$ by the ordinary flow on the interval with no intervening singular time.
  Each proper component has the finite spherical frontier and canonical
  cap/neck cover supplied by
  \ref{lem:actual-disappearing-component-frontier-saturation}.  Thus deleting each such component of
$D_{t^-}=M_{t^-}\setminus C_{t^-}$ and capping its relative boundary produces
exactly the corresponding component of $M_T$.  The relative-boundary
classification and the four explicit capping identities in
\ref{lem:local-capping-connected-sum-identities} then reconstruct the
connected-sum factors one component at a time.  Components of $D_{t^-}$ with
empty frontier are whole components and are classified by
\ref{lem:whole-disappearing-component-classification}; they are inserted as
separate standard factors.  No connected-sum conclusion is used as a premise.
Finally, $\mathbb{RP}^3\#\mathbb{RP}^3$ contributes two
$\mathbb{RP}^3=S^3/(\mathbb Z/2)$ factors by associativity.
\end{proof}

\section{Backward Reconstruction from Extinction}

\begin{theorem}[Connected-sum classification from finite extinction]
\label{thm:extinction-connected-sum-classification}
\dcref{MT:ch15:15.4,MT:intro:0.1}
\source{morgan-tian}
\uses{def:controlled-ricci-flow-with-surgery,
  def:ricci-flow-with-cutoff,
  def:smooth-connected-sum,
  def:spherical-space-form,
  def:sphere-bundles-over-circle,
  lem:two-sphere-diffeomorphism-classes,
  thm:finite-interval-surgery-bound,
  lem:single-surgery-connected-sum-reconstruction}
\notready
Let $({\mathcal M},G)$ be a fixed controlled generalized Ricci flow satisfying
Assumptions (1)--(7) and carrying the Kleiner--Lott cutoff of
\ref{def:ricci-flow-with-cutoff} on every finite restriction, with
connected closed initial slice $M_0=M$.  Assume that a regular time
$T_{\rm ext}<\infty$ has empty slice $M_{T_{\rm ext}}=\varnothing$ and that
the flow remains empty afterwards, and assume explicitly that only finitely
many singular times occur in every bounded time interval.  Then there are finite families of
spherical space forms $P_1,\ldots,P_r$ and sphere bundles
$F_1,\ldots,F_s\in\{E_+,E_-\}$ such that
\[
 M\cong P_1\#\cdots\#P_r\#F_1\#\cdots\#F_s.
\]
Here $r+s\ge1$.
If $M$ is orientable, every $F_j$ is $E_+$.
\end{theorem}

\begin{proof}
Choose $B>T_{\rm ext}$.  By \ref{thm:finite-interval-surgery-bound}, only finitely many
surgery times occur in $[0,B)$; list them as
$0<T_1<\cdots<T_k<B$.  The slice at time $B$ is empty.  Across each open
interval between consecutive surgery times, the time vector field identifies
regular slices diffeomorphically.  Starting with the empty slice after $T_k$,
apply \ref{lem:single-surgery-connected-sum-reconstruction} at a regular slice
immediately before $T_k$, transport backward through the preceding regular
interval, and repeat at $T_{k-1},\ldots,T_1$.
This finite induction expresses every earlier component, and in particular
$M$, as a finite connected sum of spherical space forms, $\mathbb{RP}^3$
factors, and $E_\pm$ factors.  Since
$\mathbb{RP}^3=S^3/(\mathbb Z/2)$, it is already a spherical space form.
Associativity of the smooth connected sum and the two sphere-gluing classes in
\ref{lem:two-sphere-diffeomorphism-classes} give the displayed form.  In the
orientable case the nonorientable bundle cannot occur, either in a
nonseparating reverse surgery or in any of the preceding connected-sum
identities.
Because the initial component $M$ is nonempty, the reconstruction contributes
at least one standard factor, so $r+s\ge1$.
\end{proof}

\begin{definition}[Thurston-geometrized closed three-manifold]
\label{def:thurston-geometrized}
\dcref{KL:app:geom,MT:ch15:15.4}
\source{kleiner-lott}
\notready
Let $M$ be a connected closed orientable smooth $3$-manifold.  We call $M$
\emph{geometrized} if there are closed connected smooth $3$-manifolds
$P_1,\ldots,P_n$, each orientable, and a diffeomorphism
$M\cong P_1\#\cdots\#P_n$ such that, for every $i$, there is a compact
codimension-zero submanifold $G_i\subset P_i$ (possibly empty or all of
$P_i$) whose nonempty boundary components are tori and for which
\begin{enumerate}
\item $G_i$ is a graph manifold: it is obtained from finitely many products
      $S^1\times P_{ij}$, with $P_{ij}$ a pair of pants, and finitely many
      products $S^1\times D^2_{ij}$ by pairing an even number of boundary
      tori by diffeomorphisms and gluing, with the gluing chosen so that
      $G_i$ is orientable;
\item the interior of $P_i\setminus G_i$ admits a complete metric of
      constant negative sectional curvature and finite volume; and
\item for every boundary torus $T\subset\partial G_i$, the inclusion
      $\pi_1(T)\to\pi_1(P_i)$ is injective.
\end{enumerate}
The cases $G_i=\varnothing$ and $G_i=P_i$ are allowed (the latter is the
closed graph-manifold case).  Equivalently, this is the finite-volume
locally-homogeneous formulation obtained by cutting along the displayed
incompressible tori.  This definition is only for connected, closed,
orientable manifolds; no orientation-double-cover convention is folded into
the term ``geometrized.''
\end{definition}

\begin{definition}[Morgan--Tian geometrized closed three-manifold]
\label{def:morgan-tian-geometrized}
\dcref{MT:ch15:15.4}
\source{morgan-tian}
\notready
A compact closed smooth $3$-manifold (possibly disconnected) is
\emph{MT-geometrized} if each connected component is a finite connected sum
of closed prime $3$-manifolds carrying one of Thurston's homogeneous
geometries.  In this predicate the two $S^2$-bundles $E_+$ and $E_-$ over
$S^1$ are admitted as the homogeneous $S^2\times\mathbb R$ factors, and a
constant-positive-curvature space form is admitted as a spherical factor.
This is the source-level predicate used in Morgan--Tian Corollary~15.4; the
more detailed orientable graph/hyperbolic package is the separate
Kleiner--Lott definition below.
\end{definition}

\begin{theorem}[Morgan--Tian Corollary 15.4, general geometrization clause]
\label{thm:morgan-tian-corollary-15-4-general}
\dcref{MT:ch15:15.4}
\source{morgan-tian}
\uses{def:ricci-flow-with-cutoff,def:morgan-tian-geometrized,
  thm:morgan-tian-proposition-15-3,thm:finite-interval-surgery-bound}
\notready
Let $({\mathcal M},G)$ be a controlled generalized Ricci flow satisfying
Assumptions (1)--(7) and carrying the Kleiner--Lott cutoff of
\ref{def:ricci-flow-with-cutoff} on $[0,T]$, with initial slice $(M,g(0))$.
If a time slice $M_T$ is
MT-geometrized, then every earlier slice $M_t$, $0\le t\le T$, is
MT-geometrized.  No orientability or connectedness hypothesis is imposed on
the individual slices beyond the assumptions of the generalized flow.
\end{theorem}

\begin{proof}
The finite-interval surgery bound gives finitely many singular times in
$[0,T]$.  Between them the time vector field gives diffeomorphisms.  At the
last surgery before $T$, Proposition~15.3 expresses the earlier slice as a
finite connected-sum assembly of the later slice, sphere bundles over $S^1$,
and positive-curvature space forms.  The latter two are homogeneous
geometrized factors, and adjoining them to a prime/geometric decomposition
preserves the defining predicate.  Induction over the finite surgery list
proves the claim.  This is exactly the backward induction in
Morgan--Tian, \texttt{surgery.tex}, lines~1600--1621.
\end{proof}

\begin{theorem}[Morgan--Tian Corollary 15.4, general empty-slice clause]
\label{thm:morgan-tian-corollary-15-4-empty-general}
\dcref{MT:ch15:15.4}
\source{morgan-tian}
\uses{thm:morgan-tian-proposition-15-3,thm:finite-interval-surgery-bound,
  def:ricci-flow-with-cutoff,def:spherical-space-form,
  def:sphere-bundles-over-circle}
\notready
Let $({\mathcal M},G)$ satisfy Assumptions (1)--(7) and carry the
Kleiner--Lott cutoff of \ref{def:ricci-flow-with-cutoff} on $[0,T]$, let the initial slice
$M_0=M$ be connected and closed, and suppose $M_T=\varnothing$ for some
$T>0$.  If only finitely many singular times occur in $[0,T]$, then
\[
 M\cong P_1\#\cdots\#P_r\#F_1\#\cdots\#F_s,
 \qquad F_j\in\{E_+,E_-\},
\]
for finitely many closed constant-positive-curvature space forms $P_i$, with
$r+s\ge1$.
\end{theorem}

\begin{proof}
Reverse the finite list of surgeries from the empty slice.  Proposition~15.3
adds only the two sphere-bundle types and positive-curvature factors, and
connected-sum associativity gives the displayed expression.  This is the
source's statement (1), \texttt{surgery.tex}, lines~1623--1632; no
orientability specialization has been inserted.
\end{proof}

\begin{lemma}[Positive space forms are geometrized]
\label{lem:positive-space-form-geometrized}
\dcref{MT:ch15:15.4}
\source{morgan-tian}
\source{kleiner-lott}
\uses{def:thurston-geometrized,def:spherical-space-form}
\notready
Every closed connected orientable three-manifold carrying a metric of constant
positive sectional curvature is geometrized in the sense of
\ref{def:thurston-geometrized}.
\end{lemma}

\begin{proof}
This is the positive-curvature-factor clause in Morgan--Tian Corollary~15.4:
the source calls these manifolds $3$-dimensional space forms and states that
they satisfy Thurston's Geometrization Conjecture
(\texttt{surgery.tex}, lines~1606--1619).  The source's homogeneous spherical
model is the $G_i=P_i$ case of the definition.
\end{proof}

\begin{lemma}[The orientable sphere bundle is geometrized]
\label{lem:orientable-sphere-bundle-geometrized}
\dcref{MT:ch15:15.4}
\source{morgan-tian}
\uses{def:thurston-geometrized,def:sphere-bundles-over-circle}
\notready
$E_+=S^2\times S^1$ is geometrized in the sense of
\ref{def:thurston-geometrized}.
\end{lemma}

\begin{proof}
Morgan--Tian Corollary~15.4 identifies each sphere bundle over $S^1$ as a
quotient of the homogeneous product $S^2\times\mathbb R$ and therefore as a
geometrized factor (\texttt{surgery.tex}, lines~1614--1619).  Explicitly,
$S^2\times S^1$ is obtained by gluing two copies of $S^1\times D^2$ along
their boundary tori, so its graph-manifold piece is the whole manifold and
the definition permits $G_i=P_i$.
\end{proof}

\begin{lemma}[Orientable geometrization is closed under reverse-surgery factors]
\label{lem:geometrization-closed-under-standard-factors}
\dcref{KL:app:geom,MT:ch15:15.3,MT:ch15:15.4}
\source{morgan-tian}
\uses{def:thurston-geometrized,def:smooth-connected-sum,
  def:spherical-space-form,def:sphere-bundles-over-circle,
  lem:two-sphere-diffeomorphism-classes,
  lem:positive-space-form-geometrized,lem:orientable-sphere-bundle-geometrized}
\notready
Let $X$ be a connected, closed, orientable geometrized $3$-manifold.  If
$Q$ is obtained from a connected-sum expression for $X$ by adjoining finitely
many orientable spherical space forms and copies of $E_+=S^2\times S^1$ and
performing finitely many connected-sum operations, then $Q$ is connected,
closed, orientable, and geometrized.
\end{lemma}

\begin{proof}
The lemmas \ref{lem:positive-space-form-geometrized} and
\ref{lem:orientable-sphere-bundle-geometrized} supply the exact source-backed
packages for every factor added by reverse surgery.

For a connected sum, retain the graph--hyperbolic package for every original
prime summand.  A spherical factor is assigned the package $G_i=P_i$, and the
factor $E_+$ is assigned its explicit graph-manifold package from
\ref{lem:orientable-sphere-bundle-geometrized}.  The registered
Kleiner--Lott statement says that a connected sum of graph manifolds is a
graph manifold and that gluing graph-manifold pieces along boundary tori
preserves that class (the graph-manifold paragraph in the geometrization
appendix).  Puncturing at a ball and performing the connected-sum gluing is
disjoint from every displayed incompressible torus, so each original
hyperbolic complement and each inclusion
$\pi_1(T)\to\pi_1(P_i)$ is unchanged.  Induction over the finite list of
additions gives the required packages, and all factors are orientable, so the
resulting connected sum is orientable.
\end{proof}

\begin{lemma}[Orientability is preserved by Ricci-flow surgery]
\label{lem:orientability-preserved-by-surgery}
\dcref{MT:ch15:15.3,MT:ch15:15.4}
\source{morgan-tian}
\uses{def:controlled-ricci-flow-with-surgery,
  lem:local-capping-connected-sum-identities}
\notready
If the initial slice of a generalized Ricci flow satisfying Assumptions
(1)--(7) is connected, closed, and orientable, then every descendant slice
before extinction is orientable.  In a reverse-surgery connected-sum
description, only the orientable bundle $E_+$ can occur.
\end{lemma}

\begin{proof}
On a regular interval the time vector field gives a diffeomorphism and hence
preserves orientability.  At a surgery time, the continuing region is obtained
by deleting necks and capping by orientable three-balls.  The local capping
identities in \ref{lem:local-capping-connected-sum-identities} show that the
reverse operation can add $E_-$ only when the resulting presurgery component
is nonorientable; an orientable component therefore receives $E_+$.
Induction over the finite surgery list proves the assertion.
\end{proof}

\begin{theorem}[Orientable specialization of Morgan--Tian Corollary 15.4]
\label{thm:morgan-tian-corollary-15-4}
\dcref{MT:ch15:15.4}
\source{morgan-tian}
\uses{def:controlled-ricci-flow-with-surgery,
  def:ricci-flow-with-cutoff,
  thm:morgan-tian-corollary-15-4-general,
  thm:morgan-tian-proposition-15-3,
  def:thurston-geometrized,
  lem:geometrization-closed-under-standard-factors}
\notready
Let $({\mathcal M},G)$ be a controlled generalized Ricci flow satisfying
Assumptions (1)--(7) and carrying the Kleiner--Lott cutoff of
\ref{def:ricci-flow-with-cutoff} on $[0,T]$, with connected, closed,
orientable initial slice $M_0$.  Assume that
there are only finitely many singular times in $[0,T]$, that every connected
component of every slice in this interval is orientable, and that every
connected component of $M_T$ is geometrized in the sense of
\ref{def:thurston-geometrized}.  Then every connected component of every
earlier slice $M_t$, $0\le t\le T$, is geometrized.  No empty-slice or
fundamental-group conclusion is part of this clause.
\end{theorem}

\begin{proof}
On every regular interval, the time vector field identifies slices by a
diffeomorphism.  At a singular time, Proposition~15.3 expresses each earlier
component as a connected sum of later components with spherical space forms
and sphere-bundle factors.  Orientability selects $E_+$, and
\ref{lem:geometrization-closed-under-standard-factors} preserves the explicit
graph--hyperbolic packages.  Induction over the finite list of singular times
proves the assertion component by component.
\end{proof}

\begin{theorem}[Orientable specialization of the Corollary 15.4 empty-slice clause]
\label{thm:morgan-tian-corollary-15-4-empty}
\dcref{MT:ch15:15.4}
\source{morgan-tian}
\uses{def:controlled-ricci-flow-with-surgery,
  def:ricci-flow-with-cutoff,
  thm:morgan-tian-corollary-15-4-empty-general,
  thm:morgan-tian-corollary-15-4,
  thm:extinction-connected-sum-classification,
  def:smooth-connected-sum,def:spherical-space-form,
  def:sphere-bundles-over-circle,lem:two-sphere-diffeomorphism-classes,
  thm:morgan-tian-proposition-15-3,lem:single-surgery-connected-sum-reconstruction,
  lem:orientability-preserved-by-surgery}
\notready
Let $({\mathcal M},G)$ be a controlled generalized Ricci flow satisfying
Assumptions (1)--(7) and carrying the Kleiner--Lott cutoff of
\ref{def:ricci-flow-with-cutoff} on $[0,T]$, with connected, closed,
orientable initial slice $M_0=M$.  Suppose
that $M_T=\varnothing$ at a time $T>0$ and that only finitely many singular
times occur in $[0,T]$.  Then
\[
 M\cong P_1\#\cdots\#P_r\#E_+^{\# s}
\]
for finitely many closed constant-positive-curvature $3$-manifolds $P_i$ and
some $s\ge0$ (with $r+s\ge1$).
\end{theorem}

\begin{proof}
The general Corollary~15.4 clause applies to the vacuously geometrized empty
terminal slice and records the geometrized status of every earlier slice.
The explicit reverse-surgery factor list is supplied by
\ref{thm:extinction-connected-sum-classification}; equivalently, it is the
orientable specialization
\ref{thm:morgan-tian-corollary-15-4} of the general clause.  Choose a regular
time after the last singular time; its slice is empty.  Across
each preceding surgery, Proposition~15.3 and
\ref{lem:single-surgery-connected-sum-reconstruction} reverse the cap and neck
operation, adding only closed positive-curvature factors and the two sphere
bundle types.  The initial and every intermediate slice are orientable, so the
twisted bundle $E_-$ cannot occur.  The finite list of singular times makes the
induction finite, and the initial component contributes at least one factor.
This is exactly the empty-slice argument in Morgan--Tian Corollary~15.4,
\texttt{surgery.tex}, lines~1623--1632.
\end{proof}

\section{Connected Sums and Fundamental Groups}

\begin{lemma}[Existence of a normalized initial metric]
\label{lem:normalized-initial-metric-exists}
\dcref{MT:ch15:p:353}
\source{morgan-tian}
\notready
Every closed smooth $3$-manifold admits a Riemannian metric which is normalized
in the sense required by the all-time controlled-surgery theorem.
\end{lemma}

\begin{proof}
Choose a locally finite smooth atlas and a subordinate partition of unity.
Pull back the Euclidean metric in every chart and take the partition-of-unity
weighted sum.  Pointwise positivity gives a smooth Riemannian metric $g$.

For $Q>0$, put $g_Q=Qg$.  Then
$|\operatorname{Rm}_{g_Q}|_{g_Q}=Q^{-1}|\operatorname{Rm}_g|_g$, so compactness
gives the required curvature bound for all sufficiently large $Q$.  Moreover,
the volume ratio of a $g$-ball of radius $\rho$ to the Euclidean ball of that
radius tends uniformly to $1$ as $\rho\downarrow0$.  A unit $g_Q$-ball
corresponds to a $g$-ball of radius $Q^{-1/2}$, with both volumes multiplied by
$Q^{3/2}$.  Taking $Q$ still larger gives the required lower bound for every
unit ball.  Thus $g_Q$ is normalized.
\end{proof}

\begin{theorem}[The free-product topology theorem]
\label{thm:free-product-connected-sum-classification}
\dcref{MT:intro:0.1}
\source{morgan-tian}
\uses{def:smooth-connected-sum,
  def:spherical-space-form,
  def:sphere-bundles-over-circle,
  lem:normalized-initial-metric-exists,
  thm:free-product-excludes-two-sided-projective-plane,
  thm:all-time-controlled-surgery-flow,
  thm:finite-time-extinction,
  thm:extinction-connected-sum-classification}
\notready
Let $M$ be a closed connected smooth $3$-manifold whose fundamental group is
a free product of finite groups and infinite cyclic groups.  Then $M$ is a
finite connected sum of spherical space forms and copies of $E_+$ and $E_-$.
If $M$ is orientable, only $E_+$ occurs.
\end{theorem}

\begin{proof}
Equip $M$ with the normalized metric from
\ref{lem:normalized-initial-metric-exists}.  The group hypothesis excludes an
embedded two-sided projective plane by
\ref{thm:free-product-excludes-two-sided-projective-plane}.  Hence
\ref{thm:all-time-controlled-surgery-flow} supplies a controlled flow with
surgery for all time.  The same group hypothesis gives finite extinction by
\ref{thm:finite-time-extinction}.  The conclusion now follows from the finite
backward reconstruction in
\ref{thm:extinction-connected-sum-classification}.
\end{proof}

\begin{theorem}[Fundamental group of a finite connected sum]
\label{thm:connected-sum-fundamental-group}
\dcref{MT:intro:0.1}
\source{morgan-tian}
\uses{def:smooth-connected-sum}
\notready
For $n\ge1$ and connected closed smooth $3$-manifolds
$M_1,\ldots,M_n$, choices of basepoints and connecting paths give an
isomorphism
\[
 \pi_1(M_1\#\cdots\#M_n)
 \cong \pi_1(M_1)*\cdots*\pi_1(M_n).
\]
Changing those choices conjugates the factor inclusions, so the isomorphism
is canonical only up to the corresponding inner automorphisms.
\end{theorem}

\begin{proof}
For two factors, cover the connected sum by the two punctured manifolds
together with a collar of the gluing sphere.  Their intersection deformation
retracts to $S^2$ and is simply connected.  Removing a ball from a connected
$3$-manifold does not change its fundamental group: apply van Kampen to the
manifold as the union of the punctured manifold and the ball, whose
intersection again retracts to $S^2$.  The Seifert--van Kampen theorem now
gives $\pi_1(M\#N)\cong\pi_1(M)*\pi_1(N)$.  Induction proves the finite formula.
\end{proof}

\begin{lemma}[A trivial free product has trivial factors]
\label{lem:trivial-free-product-factors}
\dcref{MT:intro:0.1}
\source{morgan-tian}
\notready
If a finite free product $G_1*\cdots*G_n$ is the trivial group, then every
$G_i$ is trivial.
\end{lemma}

\begin{proof}
For each $i$, the universal property of the free product gives a retraction
$G_1*\cdots*G_n\to G_i$ that is the identity on $G_i$ and sends every other
factor to the identity.  Hence the canonical map from $G_i$ into the free
product is injective.  A subgroup of the trivial group is trivial.
\end{proof}

\begin{lemma}[Fundamental group of a sphere bundle over the circle]
\label{lem:sphere-bundle-fundamental-group}
\dcref{MT:intro:0.1}
\source{morgan-tian}
\uses{def:sphere-bundles-over-circle}
\notready
Both $E_+$ and $E_-$ have fundamental group isomorphic to $\mathbb Z$.
\end{lemma}

\begin{proof}
The projection $E_\pm\to S^1$ is a fiber bundle with fiber $S^2$, hence has
the homotopy lifting property.  Its homotopy long exact sequence contains
\[
 \pi_1(S^2)\longrightarrow\pi_1(E_\pm)
 \longrightarrow\pi_1(S^1)\longrightarrow\pi_0(S^2).
\]
The outer terms are trivial, so the middle map is an isomorphism and
$\pi_1(E_\pm)\cong\mathbb Z$.
\end{proof}

\begin{lemma}[A simply connected spherical space form]
\label{lem:simply-connected-spherical-space-form}
\dcref{MT:intro:0.1}
\source{morgan-tian}
\uses{def:spherical-space-form,def:simply-connected}
\notready
If a spherical space form $P$ is simply connected, then $P$ is diffeomorphic
to $S^3$.
\end{lemma}

\begin{proof}
Write $P=S^3/\Gamma$ as in \ref{def:spherical-space-form}.  The quotient by
the free finite action is a covering map whose deck-transformation group is
$\Gamma$.  Since $S^3$ is simply connected, this is the universal cover and
$\pi_1(P)\cong\Gamma$.  If $P$ is simply connected, then $\Gamma$ is trivial,
so the quotient map itself is a diffeomorphism $S^3\to P$.
\end{proof}

\begin{lemma}[The three-sphere is the connected-sum identity]
\label{lem:three-sphere-connected-sum-identity}
\dcref{MT:ch15:15.3}
\source{morgan-tian}
\uses{def:smooth-connected-sum}
\notready
For every connected closed smooth $3$-manifold $M$,
$M\#S^3\cong M$.
\end{lemma}

\begin{proof}
The complement of the interior of a smoothly embedded ball in $S^3$ is a
ball.  Thus forming $M\#S^3$ removes a ball from $M$ and glues a ball back
along its boundary.  Since every boundary diffeomorphism extends over the
ball, the result is diffeomorphic to $M$.
\end{proof}

\section{Simple Connectivity}

\begin{theorem}[Poincare conjecture]
\label{thm:poincare-conjecture}
\dcref{MT:intro:0.2(a)}
\source{morgan-tian}
\uses{def:simply-connected,
  lem:path-connected-implies-connected,
  lem:simply-connected-three-manifold-orientable,
  lem:normalized-initial-metric-exists,
  thm:free-product-excludes-two-sided-projective-plane,
  thm:all-time-controlled-surgery-flow,
  def:ricci-flow-with-cutoff,
  thm:finite-time-extinction,
  thm:morgan-tian-corollary-15-4-empty,
  thm:free-product-connected-sum-classification,
  lem:orientability-preserved-by-surgery,
  thm:connected-sum-fundamental-group,
  lem:trivial-free-product-factors,
  lem:sphere-bundle-fundamental-group,
  lem:simply-connected-spherical-space-form,
  lem:three-sphere-connected-sum-identity}
\notready
Every compact, connected, smooth $3$-manifold without boundary which is
simply connected in the sense of \ref{def:simply-connected} is diffeomorphic
to $S^3$.
\end{theorem}

\begin{proof}
Let $M$ be such a manifold.  By
\ref{lem:path-connected-implies-connected}, simple connectedness supplies the
connectedness required by the surgery and reconstruction interfaces.  It is
orientable by
\ref{lem:simply-connected-three-manifold-orientable}.  Its trivial
fundamental group is a finite group, hence is a free product with one finite
factor.  Equip $M$ with the normalized metric from
\ref{lem:normalized-initial-metric-exists}.  The projective-plane obstruction
is excluded by \ref{thm:free-product-excludes-two-sided-projective-plane}, so
  \ref{thm:all-time-controlled-surgery-flow} produces an all-time controlled
  surgery flow with the Kleiner--Lott cutoff of
  \ref{def:ricci-flow-with-cutoff}.  It becomes extinct in finite time by
  \ref{thm:finite-time-extinction}.  The general free-product topology theorem
\ref{thm:free-product-connected-sum-classification} independently supplies
the same finite factor expression; the fixed-flow empty-slice clause below
records its source-specific reconstruction.  The empty-slice clause of
Corollary~15.4
applies directly to this same flow:
\ref{thm:morgan-tian-corollary-15-4-empty} supplies the finite standard-factor
expression under the exact connected, closed, orientable, and finite-surgery
hypotheses.  By
\ref{lem:orientability-preserved-by-surgery}, every slice in the bounded
extinction interval is orientable, so the corollary's orientability premise is
met.  Thus
\[
 M\cong P_1\#\cdots\#P_r\#E_+^{\# s}
\]
for spherical space forms $P_i$; the twisted bundle is absent by the same
orientability lemma.  The remaining elimination of factors is the elementary
group argument below, rather than a target-equivalent simply-connected
corollary imported as a premise.

By \ref{thm:connected-sum-fundamental-group}, $\pi_1(M)$ is the free product
of the groups $\pi_1(P_i)$ and $s$ copies of $\pi_1(E_+)$.  Since this free
product is trivial, \ref{lem:trivial-free-product-factors} makes every factor
trivial.  But \ref{lem:sphere-bundle-fundamental-group} gives
$\pi_1(E_+)\cong\mathbb Z$, so $s=0$.  Since $r+s\ge1$, we have $r\ge1$.
Each $P_i$ is therefore simply connected and hence diffeomorphic to $S^3$ by
\ref{lem:simply-connected-spherical-space-form}.  The connected sum is
nonempty, and repeated use of
\ref{lem:three-sphere-connected-sum-identity} now gives $M\cong S^3$.
\end{proof}
